\clearpage
\section{Analiza konkurecjnych konstrukcji}
    \paragraph{}
        Przed wykonaniem odrecznego rysunku prototypo wyselekcjonowałem 5 najbardziej zblizonych konstrukcje.
        Selekcji dokonałem za pomoca odległosci euklidesowej opisanej w literaturze \cite{analnaIza}.
        
    \subsection{Milkor 380}

        \begin{figure}[!h]
            \centering \includegraphics[width=1\linewidth]{Drony/Dron1/1.jpg}
            \caption{Milkor 380\cite{Dron1}}
        \end{figure}

        \begin{longtable}{| c | m{0,5\linewidth} | c |}
            \caption{Szczegółowe parametry Milkor 380}
            \label{table:milkor_380_specs} \\
            \hline
            \textbf{Lp} & \multicolumn{1}{c|}{\textbf{Nazwa parametru}} & \textbf{Wartość} \\ \hline \hline
            \endfirsthead
            
            \hline
            \endfoot
            
            1 & Producent & Milkor \\ \hline
            2 & Kraj produkcji & RPA \\ \hline
            3 & Dopasowanie & 93,23\% \\ \hline
            4 & Model silnika & Rotax 915 iS \\ \hline
            6 & Zasięg maksymalny  & 4000,0 km \\ \hline
            7 & Prędkość przelotowa / maksymalna & 150 / 250 km/h \\ \hline
            8 & Pułap operacyjny & 9144 m \\ \hline
            9 & Zasięg operacyjny SATA & 4000 km \\ \hline
            10 & Zasięg operacyjny LOS & 220 km \\ \hline
            11 & Długość / Wysokość / Rozpiętość & 9,00 / 2,40 / 18,60 m \\ \hline
            15 & Profil skrzydła root / tip & NACA 64-415 / NACA 64-412 \\ \hline
            16 & Masa własna & 700,0 kg \\ \hline
            17 & Masa startowa & 1300,0 kg \\ \hline
            18 & Maksymalna ładowność & 210,0 kg \\ \hline

        \end{longtable}
    
    \clearpage
    \subsection{Elbit Hermes 900}
        \begin{figure}[!h]
            \centering \includegraphics[width=1\linewidth]{Drony/Dron2/1.jpg}
            \caption{Elbit Hermes 900\cite{Dron2}}
        \end{figure}

        \begin{longtable}{| c | m{0,5\linewidth} | c |}
            \caption{Szczegółowe parametry Elbit Hermes 900}
            \label{table:hermes_900_specs} \\
            \hline
            \textbf{Lp} & \multicolumn{1}{c|}{\textbf{Nazwa parametru}} & \textbf{Wartość} \\ \hline \hline
            \endfirsthead
            
            \hline
            \endfoot
            
            1 & Producent & Elbit Systems \\ \hline
            2 & Kraj produkcji & Izrael \\ \hline
            3 & Dopasowanie & 89,99\% \\ \hline
            4 & Model silnika & Rotax 914 \\ \hline
            5 & Czas lotu & 36,0 h \\ \hline
            6 & Zasięg maksymalny & 5000,0 km \\ \hline
            7 & Prędkość przelotowa / maksymalna & 112 / 220 km/h \\ \hline
            8 & Pułap operacyjny & 9144 m \\ \hline
            9 & Zasięg operacyjny SATA & 5000 km \\ \hline
            10 & Zasięg operacyjny LOS & 250 km \\ \hline
            11 & Długość / Wysokość / Rozpiętość & 8,30 / 2,30 / 15,00 m \\ \hline
            12 & Wydłużenie płata & 15,00 \\ \hline
            13 & Obciążenie powierzchni & 78,67 kg/m$^2$ \\ \hline
            14 & Obciążenie mocy & 13,95 kg/kW \\ \hline
            15 & Profil skrzydła root / tip & Wortmann FX 63-137 \\ \hline
            16 & Masa własna & 750,0 kg \\ \hline
            17 & Masa startowa & 1180,0 kg \\ \hline
            18 & Maksymalna ładowność & 350,0 kg \\ \hline
        \end{longtable}

    \clearpage
    \subsection{Baykar Bayraktar TB2}

        \begin{figure}[!h]
            \centering \includegraphics[width=1\linewidth]{Drony/Dron3/1.jpg}
            \caption{Baykar Bayraktar TB2\cite{Dron3}}
        \end{figure}

        \begin{longtable}{| c | m{0,5\linewidth} | c |}
            \caption{Szczegółowe parametry Baykar Bayraktar TB2}
            \label{table:bayraktar_tb2_specs} \\
            \hline
            \textbf{Lp} & \multicolumn{1}{c|}{\textbf{Nazwa parametru}} & \textbf{Wartość} \\ \hline \hline
            \endfirsthead
            
            \hline
            \endfoot
            
            1 & Producent & Baykar Teknoloji \\ \hline
            2 & Kraj produkcji & Turcja \\ \hline
            3 & Dopasowanie & 89,35\% \\ \hline
            4 & Model silnika & Rotax 912 \\ \hline
            5 & Czas lotu & 27,0 h \\ \hline
            6 & Zasięg maksymalny & 4000,0 km \\ \hline
            7 & Prędkość przelotowa / maksymalna & 130 / 220 km/h \\ \hline
            8 & Pułap operacyjny & 8229 m \\ \hline
            9 & Zasięg operacyjny SATA & 4000 km \\ \hline
            10 & Zasięg operacyjny LOS & 300 km \\ \hline
            11 & Długość / Wysokość / Rozpiętość & 6,50 / 2,20 / 12,00 m \\ \hline
            12 & Wydłużenie płata & 13,71 \\ \hline
            13 & Obciążenie powierzchni & 66,67 kg/m$^2$ \\ \hline
            14 & Obciążenie mocy & 9,52 kg/kW \\ \hline
            15 & Profil skrzydła root / tip & NACA 4412 \\ \hline
            16 & Masa własna & 420,0 kg \\ \hline
            17 & Masa startowa & 700,0 kg \\ \hline
            18 & Maksymalna ładowność & 150,0 kg \\ \hline
        \end{longtable}

    \clearpage
    \subsection{Shahed 129}

        \begin{figure}[!h]
            \centering \includegraphics[width=1\linewidth]{Drony/Dron4/1.png}
            \caption{Shahed 129\cite{Dron4}}
        \end{figure}

        \begin{longtable}{| c | m{0,5\linewidth} | c |}
            \caption{SzczegółoAwe parametry Shahed 129}
            \label{table:Shahed 129} \\
            \hline
            \textbf{Lp} & \multicolumn{1}{c|}{\textbf{Nazwa parametru}} & \textbf{Wartość} \\ \hline \hline
            \endfirsthead
            
            \hline
            \endfoot
            1 & Producent & IAMIC \\ \hline
            2 & Kraj produkcji & Iran \\ \hline
            3 & Dopasowanie & 88,30\% \\ \hline
            4 & Model silnika & Rotax 914 \\ \hline
            5 & Czas lotu & 24,0 h \\ \hline
            6 & Zasięg maksymalny & 1700,0 km \\ \hline
            7 & Prędkość przelotowa / maksymalna & 150 / 175 km/h \\ \hline
            8 & Pułap operacyjny & 7300 m \\ \hline
            9 & Zasięg operacyjny SATA & 1700 km \\ \hline
            10 & Zasięg operacyjny LOS & 200 km \\ \hline
            11 & Długość / Wysokość / Rozpiętość & 8,00 / 3,10 / 16,00 m \\ \hline
            12 & Wydłużenie płata & 14,88 \\ \hline
            13 & Obciążenie powierzchni & 63,95 kg/m$^2$ \\ \hline
            14 & Obciążenie mocy & 13,01 kg/kW \\ \hline
            15 & Masa własna & 700,0 kg \\ \hline
            16 & Masa startowa & 1100,0 kg \\ \hline
            17 & Maksymalna ładowność & 400,0 kg \\ \hline
        \end{longtable}

    \clearpage
    \subsection{Wing Loong I}

        \begin{figure}[!h]
            \centering \includegraphics[width=1\linewidth]{Drony/Dron5/1.jpg}
            \caption{Shahed 129\cite{Dron5}}
        \end{figure}

        \begin{longtable}{| c | m{0,5\linewidth} | c |}
            \caption{SzczegółoAwe parametry Shahed 129}
            \label{table:Wing Loong I} \\
            \hline
            \textbf{Lp} & \multicolumn{1}{c|}{\textbf{Nazwa parametru}} & \textbf{Wartość} \\ \hline \hline
            \endfirsthead
            
            \hline
            \endfoot
            1 & Producent & CAIG \\ \hline
            2 & Kraj produkcji & Chiny \\ \hline
            3 & Dopasowanie & 87,77\% \\ \hline
            4 & Model silnika & Rotax 914 \\ \hline
            5 & Czas lotu & 20,0 h \\ \hline
            6 & Zasięg maksymalny & 4000,0 km \\ \hline
            7 & Prędkość przelotowa / maksymalna & 150 / 280 km/h \\ \hline
            8 & Pułap operacyjny & 7500 m \\ \hline
            9 & Zasięg operacyjny SATA & 4000 km \\ \hline
            10 & Zasięg operacyjny LOS & 200 km \\ \hline
            11 & Długość / Wysokość / Rozpiętość & 9,05 / 2,70 / 14,00 m \\ \hline
            12 & Wydłużenie płata & 12,17 \\ \hline
            13 & Obciążenie powierzchni & 68,32 kg/m$^2$ \\ \hline
            14 & Obciążenie mocy & 13,01 kg/kW \\ \hline
            15 & Profil skrzydła root / tip & Drela GW-19 / Drela GW-27 \\ \hline
            16 & Masa własna & 600,0 kg \\ \hline
            17 & Masa startowa & 1100,0 kg \\ \hline
            18 & Maksymalna ładowność & 200,0 kg \\ \hline
        \end{longtable}

    \clearpage
    \subsection{Model prototypu}

    \begin{figure}[!h]
        \centering \includegraphics[width=1\linewidth]{Model/model.pdf}
        \caption{Model prototypu}
    \end{figure}

    \paragraph{}
        W strukturze projektowanego BSP zdecydowano się na zastosowanie układu centropłata, 
        Wybór ten podyktowany jest przede wszystkim pozwoli osiągnąć optymalne osiągi aerodynamiczne oraz wysokią manewrowości
        oraz wyposazony w chowane podwozie o budowie,  
        Konfiguracja ta wiąże się z ograniczeniem przestrzeni użytecznej wewnątrz kadłuba,
        Omawiany dron nie jest przeznaczony do przewozu wielkogabarytowego ładunku płatnego, 
        lecz do obserwacji, 
        Wyposażenie stanowi jedynie zintegrowana głowica optoelektroniczna,
    
    \paragraph{}
        Niemniej jednak, mając na uwadze konieczność optymalizacji czasu trwania procesu projektowego oraz dążenie do minimalizacji ryzyka błędu, 
        zdecydowano się na implementację usterzenia klasycznego, 
    
    \paragraph{}
        W projekcie zastosowano układ ze śmigłem pchającym z tyłu kadłuba, 
        Decyzja ta wynika z konieczności zapewnienia niezakłóconego pola widzenia dla głowicy optoelektronicznej, 
        Przeniesienie napędu na ogon uwalnia przestrzeń na dziobie drona, 
        eliminując ryzyko wchodzenia łopat śmigła w kadr i powstawania zakłóceń obrazu, 
        Dodatkowo chroni to delikatną optykę przed zanieczyszczeniami i uszkodzeniami podczas lądowania,